\chapter{乘积测度、富比尼、变量替换}
\label{ch:b3:product}

借助两个定理，一维勒贝格理论化为多维
微积分。\emph{托内利--富比尼}断言，
乘积上的积分即为累次积分——切片
在任一顺序下皆合法，前提是可以实际
检验的假设。\emph{变量替换公式}把
积分沿 $\mathcal C^1$ 微分\omterm{def:b3:topology:continuity}{同胚}输运，以雅可比行列式
作为体积的兑换率，并完整证明之，
先从线性情形出发，在那里它解释了
行列式\emph{究竟是什么}。应用随之展开：层饼公式、
卷积、\omterm{ex:b3:product:polar}{极坐标}、$n$ 维球的
体积——以及周末问题中的斯特林公式，
并给出诚实的误差分析。

\section{乘积 $\sigma$-代数与乘积测度}

\begin{definition}\label{def:b3:product:sigma}
设 $(X, \mathcal A)$、$(Y, \mathcal B)$ 为可测空间，则
\emph{乘积 $\sigma$-代数}\index{乘积
$\sigma$-代数} $\mathcal A \otimes \mathcal B$ 是 $X \times
Y$ 上由\emph{矩形} $A \times B$（$A \in
\mathcal A$，$B \in \mathcal B$）生成——它是一个 $\pi$-系。对 $E
\subseteq X\times Y$ 及 $x \in X$，\emph{截口}为 $E_x
= \{y : (x,y) \in E\}$；对乘积上的函数 $f$，
记 $f_x = f(x, \cdot)$。
\end{definition}

\begin{proposition}\label{prop:b3:product:sections}
(a) 若 $E \in \mathcal A\otimes\mathcal B$，则每个截口 $E_x
\in \mathcal B$（对另一个变量同理）；若 $f$ 关于 $\mathcal
A\otimes\mathcal B$ 可测，则每个 $f_x$ 关于 $\mathcal
B$ 可测。
(b) $\mathcal B(\R^m)\otimes\mathcal B(\R^n) = \mathcal
B(\R^{m+n})$.
\end{proposition}

\begin{proof}
(a) 良好集：$\{E : E_x \in \mathcal B\ \forall x\}$ 是一个
$\sigma$-代数（截口与取补、
可数并运算可交换），且包含所有矩形。对 $f$：
$(f_x)^{-1}(B) = (f^{-1}(B))_x$。
(b) （$\subseteq$）博雷尔集构成的矩形：只需说明
开$\times$开方体在 $\R^{m+n}$ 中是博雷尔集（它们是\omterm{def:b3:topology:topology}{开集}），
且一般的博雷尔矩形都是其极限——
又是良好集：$\{A : A\times\R^n \in \mathcal B(\R^{m+n})\}$ 是
$\sigma$-代数，包含所有\omterm{def:b3:topology:topology}{开集}；两个这样的集合取交即可。
（$\supseteq$）$\R^{m+n}$ 中每个\omterm{def:b3:topology:topology}{开集}都是可数个
有理开方体 $U\times V$ 的并，
因而属于乘积 $\sigma$-代数。
\end{proof}

\begin{theorem}[乘积测度]\label{thm:b3:product:existence}
设 $(X, \mathcal A, \mu)$ 与 $(Y, \mathcal B, \nu)$ 为
\emph{$\sigma$-有限}\omterm{def:b3:measure:measure}{测度空间}。对每个 $E \in \mathcal
A\otimes\mathcal B$，函数 $x \mapsto \nu(E_x)$ 是
可测的，并且\index{乘积测度}
\[
(\mu\otimes\nu)(E) = \int_X \nu(E_x)\,\dd\mu(x)
\]
它定义了 $\mathcal A\otimes\mathcal B$ 上唯一的\omterm{def:b3:measure:measure}{测度}，
满足 $(\mu\otimes\nu)(A\times B) = \mu(A)\nu(B)$。它
是 $\sigma$-有限的，并且对称：
对 $x$-截口关于 $\nu$ 积分，所得\omterm{def:b3:measure:measure}{测度}相同。
\end{theorem}

\begin{proof}
\emph{$x \mapsto \nu(E_x)$ 的可测性。}首先设 $\nu$
为有限\omterm{def:b3:measure:measure}{测度}。记 $\mathcal D$ 为使映射可测的 $E$ 构成的类，
它包含矩形（$\nu((A\times B)_x) =
\nu(B)\mathbf 1_A(x)$），并且是 $\lambda$-系：对 $E
\subseteq F$ 属于 $\mathcal D$，$\nu((F\setminus E)_x) =
\nu(F_x) - \nu(E_x)$（由有限性）；对 $E_n \uparrow E$，
$\nu((E_n)_x) \uparrow \nu(E_x)$（由下连续性），并且
\omterm{def:b3:lebesgue:measurable}{可测函数}的单调极限仍可测。
矩形构成一个 $\pi$-系。由狄恩金
（\cref{thm:b3:measure:dynkin}）得 $\mathcal D = \mathcal
A\otimes\mathcal B$。若 $\nu$ 为 $\sigma$-有限的，记 $Y =
\bigcup Y_k$，$Y_k \uparrow$，$\nu(Y_k) < \infty$：$\nu(E_x) =
\lim_k\nu_k(E_x)$，其中 $\nu_k = \nu(\cdot\cap Y_k)$ 有限。

\emph{\omterm{def:b3:measure:measure}{测度}。}$\sigma$-可加性：$E \mapsto
\int\nu(E_x)\dd\mu$ 由
\cref{cor:b3:lebesgue:additivity}（不相交集的截口互不相交）推出。在矩形上，它给出
$\mu(A)\nu(B)$。
\emph{唯一性}：两个候选\omterm{def:b3:measure:measure}{测度}在矩形构成的 $\pi$-系上一致；
由 $\sigma$-有限性，取有限\omterm{def:b3:measure:measure}{测度}的矩形 $X_k\times
Y_k \uparrow X\times Y$（\omterm{def:b3:measure:measure}{测度}有限）：
由\cref{thm:b3:measure:uniqueness}可知两者相同。对称性：
另一顺序的构造同样是\omterm{def:b3:measure:measure}{测度}，且在矩形上一致——
因而唯一，从而二者相同。
\end{proof}

\begin{definition}\label{def:b3:product:lebesgue}
\emph{$\R^d$ 上的勒贝格测度}\index{
$\R^d$ 上的勒贝格测度}是 $\lambda_d = \lambda\otimes\cdots\otimes\lambda$
（$d$ 个因子；构造的结合性在方体上验证，
并由唯一性推广）。它是唯一的博雷尔\omterm{def:b3:measure:measure}{测度}，使每个方体
$\prod\intoc{a_i}{b_i}$ 的体积为
$\prod(b_i - a_i)$；它是平移不变的（平移后的\omterm{def:b3:measure:measure}{测度}
在方体上一致），$\sigma$-有限，并在
卡拉泰奥多里\omterm{thm:b3:complete:completion}{完备化}后成为完备\omterm{def:b3:measure:measure}{测度}——仍以 $\lambda_d$ 记\omterm{thm:b3:complete:completion}{完备化}后的\omterm{def:b3:measure:measure}{测度}，
并据此积分。
\end{definition}

\section{托内利与富比尼}

\begin{theorem}[托内利]\label{thm:b3:product:tonelli}
设 $\mu, \nu$ 为 $\sigma$-有限\omterm{def:b3:measure:measure}{测度}，$f \colon X\times Y \to [0,
+\infty]$ 为\omterm{def:b3:lebesgue:measurable}{可测函数}。则 $x \mapsto \int_Y f_x\,\dd\nu$ 是
可测的，并且\index{托内利定理}
\[
\int_{X\times Y}f\,\dd(\mu\otimes\nu)
= \int_X\Bigl(\int_Y f(x,y)\,\dd\nu(y)\Bigr)\dd\mu(x)
= \int_Y\Bigl(\int_X f(x,y)\,\dd\mu(x)\Bigr)\dd\nu(y).
\]
\end{theorem}

\begin{proof}
标准论证。对 $f = \mathbf 1_E$，此式
即为\cref{thm:b3:product:existence}（及其对称形式）。由线性性，
它对\omterm{def:b3:lebesgue:simple}{简单函数} $f \geq 0$ 成立。对一般的 $f \geq
0$：取\omterm{def:b3:lebesgue:simple}{简单函数}列 $s_n \nearrow f$
（见\cref{thm:b3:lebesgue:approximation}）；则 $\int_Y(s_n)_x
\dd\nu \nearrow \int_Y f_x\dd\nu$ 对每个 $x$ 成立（由 $Y$ 上的单调收敛定理），
故左端由 $X$ 上的单调收敛定理收敛，同时
$\int s_n\,\dd(\mu\otimes\nu) \nearrow \int f$ 由
\omterm{thm:b3:product:existence}{乘积测度}上的单调收敛定理得出。
\end{proof}

\begin{theorem}[富比尼]\label{thm:b3:product:fubini}
$\mu, \nu$ 为 $\sigma$-有限，$f \in L^1(\mu\otimes\nu)$。则
对 $\mu$-几乎处处的 $x$，截面 $f_x$ 是 $\nu$-可积的，
几乎处处有定义的函数 $x \mapsto \int f_x\dd\nu$ 是
\omterm{def:b3:lebesgue:l1}{可积}的，且两个累次积分都等于
$\int f\,\dd(\mu\otimes\nu)$。\index{富比尼定理}
\end{theorem}

\begin{proof}
对 $\abs f$ 用托内利定理，可知 $\varphi(x) = \int\abs{f_x}
\dd\nu$ 的积分有限，故几乎处处有限，即 $f_x \in
L^1(\nu)$ 对几乎处处的 $x$ 成立。取 $f = f^+ - f^-$（实值情形；
复值情形按分量处理）。由托内利定理，$f^\pm$ 的每个累次积分
都等于 $\int f^\pm\dd(\mu\otimes\nu) <
\infty$，几乎处处定义的差的积分等于
两者之差。另一顺序同理。
\end{proof}

\begin{method}\label{met:b3:product:fubini}
交换两个积分（或一个积分与一个和，或两个和式）的次序：若被积函数
\emph{非负}，则可随意交换（托内利定理）。否则，
先以较易估计的次序对 $\abs f$ 用托内利定理；
若结果有限，
则由富比尼定理保证交换合法。切勿略去
$\abs f$ 的检验：\cref{exo:b3:product:4}中的被积函数含两个累次积分，
二者\emph{值不同}。
\end{method}

\begin{proposition}[层蛋糕公式]\label{prop:b3:product:layercake}
设 $f \geq 0$ 为 $(X, \mathcal A, \mu)$ 上的\omterm{def:b3:lebesgue:measurable}{可测函数}，
且\omterm{def:b3:measure:measure}{测度}为 $\sigma$-有限：\index{层蛋糕公式}
\[
\int_X f\,\dd\mu = \int_0^{+\infty}\mu(\{f > t\})\,\dd t,
\qquad
\int_X f^p\,\dd\mu = p\int_0^{+\infty}t^{p-1}\mu(\{f >
t\})\,\dd t \quad (p \geq 1).
\]
\end{proposition}

\begin{proof}
对 $\mathbf 1_{\{(x,t) : 0 < t < f(x)\}}$ 在 $X
\times \intoo0{+\infty}$ 上用托内利定理（可测性：它是 $\{(x,t): f(x) - t
> 0\}\cap\{t > 0\}$，即\omterm{def:b3:lebesgue:measurable}{可测函数}的博雷尔型组合，
$(x,t)\mapsto f(x) - t$ 可测）。先对 $t$ 积分，得 $\int f\,\dd\mu$；
先对 $x$ 积分，则为 $\int_0^\infty\mu(f >
t)\dd t$。对 $f^p$，取 $t = s^p$ 代入 $\int\mu(f^p > t)
\dd t$，即对 $f^p$ 运用第一式，并
在一维积分中换元（利用 $\{f^p > s^p\} =
\{f > s\}$）。
\end{proof}

\begin{theorem}[$L^1$ 上的卷积]\label{thm:b3:product:convolution}
设 $f, g \in L^1(\R^d, \lambda_d)$，积分
\[
(f * g)(x) = \int_{\R^d} f(x - y)\,g(y)\,\dd y
\]
对几乎处处的 $x$ 绝对收敛，定义 $f * g \in
L^1(\R^d)$，且 $\norm{f*g}_1 \leq \norm f_1\norm g_1$，并且
$*$ 满足交换律与结合律。\index{卷积}
\end{theorem}

\begin{proof}
$(x, y) \mapsto f(x-y)g(y)$ 可测（$(x,y)\mapsto x -
y$ 连续，经复合与乘积即得）。托内利定理给出：
\[
\int\!\!\int \abs{f(x-y)}\abs{g(y)}\,\dd y\,\dd x
= \int\abs{g(y)}\Bigl(\int\abs{f(x - y)}\dd x\Bigr)\dd y
= \norm f_1\norm g_1 < \infty
\]
（内层积分中利用 $\lambda_d$ 的平移不变性）。
于是二重积分有限；由富比尼定理，几乎处处绝对
收敛，并得范数估计 $\norm{f*g}_1 \leq \norm
f_1\norm g_1$。交换律：作代换 $y \mapsto x - y$
（平移不变性与反射不变性——反射不变性在方体上成立，
故由唯一性处处成立）。
结合律：对三重积分用托内利--富比尼定理。
\end{proof}

\section{变量替换}

\begin{theorem}[线性变量替换]
\label{thm:b3:product:linearchange}
设 $T \in GL_d(\R)$，$A \in \mathcal B(\R^d)$：
$\lambda_d(T(A)) = \abs{\det T}\,\lambda_d(A)$；因此
$\int f(y)\dd y = \abs{\det T}\int f(Tx)\,\dd x$，对 $f \geq 0$
或\omterm{def:b3:lebesgue:l1}{可积}时成立。\index{线性变量替换}
\end{theorem}

\begin{proof}
记 $\mu_T(A) = \lambda_d(T(A))$，它是博雷尔\omterm{def:b3:measure:measure}{测度}
（\omterm{def:b3:topology:continuity}{同胚}保持博雷尔集，
\cref{pb:b3:measure:1}），平移不变
（$T(A + x) = T(A) + Tx$），在单位方体上有限。由
\omterm{def:b3:measure:lebesgueouter}{勒贝格测度}的刻画（
\cref{exo:b3:measure:6}，其证明在
$\R^d$ 中用二进方体照搬即可），$\mu_T = c(T)\lambda_d$，其中 $c(T)
= \lambda_d(T(\intco01^d))$。映射 $T \mapsto c(T)$ 是
乘性的（$c(ST) = c(S)c(T)$，由复合可得），故
只需计算 $c$ 在 $GL_d$ 的生成元上的值：初等
矩阵。对角矩阵 $\operatorname{diag}(a, 1, \dots, 1)$：
将单位立方体映为体积 $\abs a$ 的方体：$c = \abs a =
\abs\det$。坐标的对换：置换立方体，$c
= 1 = \abs\det$。切变 $T(x) = x + \alpha
x_2e_1$：单位立方体的像为一剪切棱柱；由
托内利定理，其\omterm{def:b3:measure:measure}{测度}为 $\int\lambda_1(\text{截面})\dd x_2
\cdots \dd x_d = 1$，每个 $x_1$ 截面都是
长为 $1$ 的区间：$c = 1 = \abs{\det}$。每个可逆矩阵都是
这些矩阵的乘积（高斯消元），且 $c$ 与
$\abs\det$ 都是乘性的，于是 $c(T) = \abs{\det T}$。
积分公式由标准论证得出（依次对示性函数、
\omterm{def:b3:lebesgue:simple}{简单函数}及单调收敛定理）。
\end{proof}

\begin{theorem}[变量替换]\label{thm:b3:product:changeofvar}
设 $U, V \subseteq \R^d$ 为\omterm{def:b3:topology:topology}{开集}，$\Phi \colon U \to V$ 为
$\mathcal C^1$ 类微分\omterm{def:b3:topology:continuity}{同胚}。对每个\omterm{def:b3:lebesgue:measurable}{可测函数} $f \colon V
\to [0, +\infty]$（或 $f \in L^1(V)$）：\index{变量
替换公式}
\[
\int_V f(y)\,\dd y = \int_U
f\bigl(\Phi(x)\bigr)\,\abs{\det D\Phi(x)}\,\dd x .
\]
\end{theorem}

\begin{proof}
记 $J(x) = \abs{\det D\Phi(x)}$。证明的核心是
下面的不等式：
\[
\lambda_d\bigl(\Phi(A)\bigr) \leq \int_A J\,\dd\lambda_d
\qquad\text{对每个博雷尔集 } A \subseteq U;
\tag{$*$}
\]
下面第 4 步把 $(*)$——分别用于 $\Phi$ 与
$\Phi^{-1}$——，推出定理中的等式。注意
$\Phi^{-1}$ 本身也是 $\mathcal C^1$ 类微分\omterm{def:b3:topology:continuity}{同胚}，其
雅可比行列式为 $\abs{\det D\Phi^{-1}(y)} = J(\Phi^{-1}y)^{-1}$
（对 $\Phi\circ\Phi^{-1} = \mathrm{id}$ 用链式法则）。

\emph{第 1 步：带畸变因子的方体上的 $(*)$。}取一个
闭方体 $Q \subseteq U$，中心为 $x_0$，棱长为 $2r$
（上确界范数意义下的球）。断言：对任意 $\varepsilon > 0$，若 $\Phi$
在 $Q$ 上可微，且
$\norm{D\Phi(x) - D\Phi(x_0)} \leq \varepsilon$ 在 $Q$ 上
（对上确界范数取\omterm{def:b3:banach:operator}{算子范数}），则
\[
\Phi(Q) \subseteq \Phi(x_0) + D\Phi(x_0)\Bigl(\,\bigl(1 +
\varepsilon\norm{D\Phi(x_0)^{-1}}\bigr)\,(Q - x_0)\Bigr),
\]
这是因为对 $x \in Q$，将中值不等式用于
$\Phi(x) - \Phi(x_0) - D\Phi(x_0)(x - x_0)$，得
$\norm{\Phi(x) - \Phi(x_0) - D\Phi(x_0)(x-x_0)}_\infty \leq
\varepsilon\norm{x - x_0}_\infty \leq \varepsilon r$，且
$D\Phi(x_0)^{-1}$ 把这一偏差化为
对方体作 $\varepsilon\norm{D\Phi(x_0)^{-1}}\,r$ 的放大。
由\cref{thm:b3:product:linearchange}，
\[
\lambda_d(\Phi(Q)) \leq \abs{\det D\Phi(x_0)}\,
\bigl(1 + \varepsilon\,C\bigr)^{d}\,\lambda_d(Q),
\qquad C = \sup_{Q}\norm{D\Phi(\cdot)^{-1}} .
\]

\emph{第 2 步：紧方体上的 $(*)$，由剖分得到。}设 $Q
\subseteq U$ 为紧方体，$\varepsilon > 0$。在 $Q$ 上，
$D\Phi$ 一致连续，且 $\norm{D\Phi^{-1}}$
有界（由紧性）；将 $Q$ 剖分为 $2^{kd}$ 个子方体
$Q_i$，使每个子方体上 $D\Phi$ 的振幅
$\leq \varepsilon$。对每个子方体（中心为 $x_i$）用第 1 步：
\[
\lambda_d(\Phi(Q)) \leq \sum_i\lambda_d(\Phi(Q_i))
\leq (1 + C\varepsilon)^d \sum_i \abs{\det
D\Phi(x_i)}\,\lambda_d(Q_i)
\leq (1 + C\varepsilon)^d\Bigl(\int_Q J + \varepsilon'\Bigr),
\]
末一步的依据是 $\sum_i\abs{\det D\Phi(x_i)}\mathbf
1_{Q_i} \to J$ 在 $Q$ 上一致收敛（由 $\det D\Phi$ 连续）
——黎曼和比较。令 $\varepsilon \to 0$：$(*)$
对紧方体成立。

\emph{第 3 步：$(*)$ 对一切博雷尔集 $A$ 成立。}集函数 $A
\mapsto \lambda_d(\Phi(A))$，在 $U$ 的博雷尔子集上，是一个
\omterm{def:b3:measure:measure}{测度}（$\Phi$ 是到 $V$ 上的双射，保持博雷尔集与
可数不交性），$A \mapsto \int_AJ$ 亦然。
$U$ 的每个开子集都是可数个
几乎不交的二进紧方体之并（标准二进分解：
取含于该\omterm{def:b3:topology:topology}{开集}的极大二进方体），
两个\omterm{def:b3:measure:measure}{测度}在它们之上都可加（方体的边界
为 $\lambda_d$ 零测集，其 $\Phi$ 像也是零测集，
由第 2 步对各面的薄方体覆盖得到）：
$(*)$ 由方体推广到\omterm{def:b3:topology:topology}{开集}。一般的博雷尔集 $A$：
取 $U$ 的紧穷竭 $K_m \uparrow U$，使 $K_m \subseteq
\mathring K_{m+1}$，固定 $m$；在 $\mathring K_{m+1}$ 上，$J$
有界，界为 $M_m$。由 $\lambda_d$ 的外正则性
（证明同\cref{thm:b3:measure:regularity}，以箱体代替），
取\omterm{def:b3:topology:topology}{开集} $O_n$，使 $A \cap K_m \subseteq O_n
\subseteq \mathring K_{m+1}$，且 $\lambda_d\bigl(O_n \setminus
(A\cap K_m)\bigr) \to 0$。于是
\[
\lambda_d\bigl(\Phi(A\cap K_m)\bigr) \leq
\lambda_d\bigl(\Phi(O_n)\bigr) \leq \int_{O_n}J
\leq \int_{A\cap K_m}J + M_m\,\lambda_d\bigl(O_n\setminus(A\cap
K_m)\bigr) \xrightarrow[n\to\infty]{} \int_{A\cap K_m}J .
\]
令 $m \to \infty$：左边用自下连续性，右边用
单调收敛定理。由此得 $(*)$。

\emph{第 4 步：等式与积分公式。}先
将 $(*)$ 由集合推广到积分：对每个\omterm{def:b3:lebesgue:measurable}{可测函数} $g
\geq 0$（定义于 $V$ 上），
\[
\int_V g(y)\,\dd y \leq \int_U g(\Phi(x))\,J(x)\,\dd x .
\tag{$**$}
\]
事实上，当 $g = \mathbf 1_B$ 时，它即 $(*)$ 中取 $A =
\Phi^{-1}(B)$ 的情形；由线性性推广到\omterm{def:b3:lebesgue:simple}{简单函数} $g$，再由
单调收敛定理推广到一切 $g \geq 0$（标准方法）。现对 $(**)$ 应用两次：
先用于 $g$，再——对微分\omterm{def:b3:topology:continuity}{同胚} $\Phi^{-1}$——
用于函数 $x \mapsto g(\Phi(x))J(x)$：
\[
\int_Vg \leq \int_U g(\Phi(x))J(x)\dd x
\leq \int_V g(y)\,J(\Phi^{-1}y)\,\abs{\det
D\Phi^{-1}(y)}\,\dd y = \int_V g ,
\]
由于 $J(\Phi^{-1}y)\abs{\det D\Phi^{-1}(y)} = \abs{\det\bigl(
D\Phi(\Phi^{-1}y)\,D\Phi^{-1}(y)\bigr)} = 1$（链式法则用于
$\Phi\circ\Phi^{-1} = \mathrm{id}$）。所有不等式均
取等号：公式对 $g \geq 0$ 成立，对 $L^1$
函数由分解得出。
\end{proof}

\begin{example}[极坐标；再谈高斯积分]
\label{ex:b3:product:polar}
$\Phi(r, \theta) = (r\cos\theta, r\sin\theta)$ 是 $\mathcal
C^1$ 类微分\omterm{def:b3:topology:continuity}{同胚}，从 $\intoo0{+\infty}\times\intoo0{2\pi}$
满射至 $\R^2$ 去掉一条半直线（零测集）后的集合，且 $\det D\Phi =
r$：\index{极坐标}
\[
\int_{\R^2}f(x, y)\,\dd x\,\dd y
= \int_0^{2\pi}\!\!\int_0^{+\infty}
f(r\cos\theta, r\sin\theta)\,r\,\dd r\,\dd\theta .
\]
设 $f = \eu^{-x^2-y^2}$，由托内利定理及上式可得
\[
G^2 = \Bigl(\int_\R \eu^{-x^2}\dd x\Bigr)^2
= \int_{\R^2}\eu^{-x^2-y^2}
= 2\pi\int_0^\infty r\eu^{-r^2}\dd r = \pi:
\]
这是 $G = \sqrt\pi$ 的经典两行证明，现在已
完全严格（可与
\cref{pb:b3:lebesgue:1} 中的参数证明比较）。
\end{example}

\begin{theorem}[单位球的体积]
\label{thm:b3:product:ballvolume}
设 $v_d = \lambda_d(B(0,1))$ 为 $\R^d$ 中单位球的体积，则\index{单位球的
体积}
\[
v_d = \frac{\pi^{d/2}}{\Gamma\bigl(\frac d2 + 1\bigr)} :
\qquad
v_1 = 2,\quad v_2 = \pi,\quad v_3 = \tfrac{4\pi}3,\quad
v_4 = \tfrac{\pi^2}2,\ \dots
\]
\end{theorem}

\begin{proof}
以两种方式计算 $I = \int_{\R^d}\eu^{-\norm x_2^2}\dd\lambda_d$。
由托内利定理，它分解为 $I = G^d = \pi^{d/2}$。利用
层饼公式（\cref{prop:b3:product:layercake}），取 $f = \eu^{-
\norm x^2}$，其水平集为球：$\{f > t\} =
B\bigl(0, \sqrt{-\ln t}\bigr)$ 对 $0 < t < 1$，其\omterm{def:b3:measure:measure}{测度}为
$v_d(-\ln t)^{d/2}$（以 $\rho$ 作伸缩时 $\lambda_d$ 乘以
$\rho^d$：见\cref{thm:b3:product:linearchange}），于是
\[
I = \int_0^1 v_d\,(-\ln t)^{d/2}\,\dd t
\overset{t = \eu^{-s}}{=}
v_d\int_0^\infty s^{d/2}\eu^{-s}\,\dd s
= v_d\,\Gamma\Bigl(\frac d2 + 1\Bigr).
\]
两式相等。（取值：$\Gamma(\frac32) = \frac{\sqrt\pi}2$，
$\Gamma(2) = 1$ 等）。注意 $v_d \to 0$（$d \to \infty$ 时）——
周末问题借助斯特林公式给出其收敛速度。
\end{proof}

\section{练习}

\begin{exercise}[$\star$]\label{exo:b3:product:1}
设 $\mu$ 为 $(\intcc01, \mathcal
B(\intcc01))$ 上的计数\omterm{def:b3:measure:measure}{测度}（非 $\sigma$-有限），$\lambda$ 为\omterm{def:b3:measure:lebesgueouter}{勒贝格
测度}，记 $\Delta = \{(x,x)\}$ 为 $\intcc01^2$ 中的
对角线。证明 $\Delta$ 可测，并计算 $\mathbf 1_\Delta$ 分别关于
$\lambda$ 与 $\mu$ 的两个累次积分，二者不相等。
\cref{thm:b3:product:tonelli} 的哪一条
假设不成立？
\end{exercise}

\begin{exercise}[$\star$]\label{exo:b3:product:2}
证明交换积分次序的合理性，并重新推导狄利克雷积分：对
$A > 0$，
\[
\int_0^A\frac{\sin x}x\,\dd x
= \int_0^A\!\!\int_0^{+\infty}\eu^{-xy}\sin x\,\dd y\,\dd x
= \int_0^{+\infty}\!\!\int_0^A \eu^{-xy}\sin x\,\dd x\,\dd y,
\]
求内层积分的闭形式，并令 $A \to
+\infty$（对 $y$ 积分作控制），得到
$\int_0^\infty\frac{\sin x}x\dd x = \frac\pi2$。
\end{exercise}

\begin{exercise}[$\star\star$]\label{exo:b3:product:3}
(a) 证明：若 $f \geq 0$ 可测且 $\mu$ 有限，则
$\sum_{n\geq1}\mu(\{f \geq n\}) \leq \int f\,\dd\mu \leq
\mu(X) + \sum_{n\geq1}\mu(\{f\geq n\})$：\omterm{def:b3:lebesgue:l1}{可积}性即
尾\omterm{def:b3:measure:measure}{测度}的可和性。
(b) 由此推得 $f \in L^1(\mu)$（$\mu$ 有限）当且仅当
$\sum_n\mu(\abs f \geq n) < \infty$。
\end{exercise}

\begin{exercise}[$\star\star$]\label{exo:b3:product:4}
设 $f(x, y) = \dfrac{x^2 - y^2}{(x^2 + y^2)^2}$ 在
$\intoo01^2$ 上，证明
\[
\int_0^1\!\!\int_0^1 f\,\dd y\,\dd x = \frac\pi4,
\qquad
\int_0^1\!\!\int_0^1 f\,\dd x\,\dd y = -\frac\pi4
\]
\emph{（注意 $f = \partial_y\bigl(\frac{y}{x^2+y^2}\bigr)$）}，
并直接验证 $\int\!\int\abs f = +\infty$：富比尼
的\omterm{def:b3:lebesgue:l1}{可积}性假设绝非装饰。
\end{exercise}

\begin{exercise}[$\star\star$]\label{exo:b3:product:5}
(a) 显式计算 $\mathbf 1_{\intcc01} * \mathbf 1_{\intcc01}$
（帐篷函数），并给出 $(\mathbf 1 * \mathbf 1 *
\mathbf 1)$ 的一般形状。
(b) 证明 $\operatorname{supp}(f * g) \subseteq
\overline{\operatorname{supp}f + \operatorname{supp}g}$。
(c) 证明：若 $f \in L^1$，$g$ 有界且连续，
则 $f * g$ 连续。\emph{（由控制收敛定理，借助
有界函数 $g$ 的平移连续性。）}
\end{exercise}

\begin{exercise}[$\star\star$]\label{exo:b3:product:6}
(a) 证明单形 $\Delta_d = \{x \in \intco0\infty^d
: x_1 + \dots + x_d \leq 1\}$ 的体积为 $\frac1{d!}$
\emph{（归纳法与富比尼定理）}。
(b) 恢复 $v_2 = \pi$、$v_3 = \frac{4\pi}3$（见
\cref{thm:b3:product:ballvolume}），并证明
$\lambda_d(\text{椭球，半轴为 } a_i) =
v_d\prod a_i$。
\end{exercise}

\begin{exercise}[$\star\star$]\label{exo:b3:product:7}
对哪些 $s > 0$，下列积分有限？请用\omterm{ex:b3:product:polar}{极
坐标}说明理由：
\[
\int_{B(0,1)\subseteq\R^2}\frac{\dd x\,\dd y}{(x^2 +
y^2)^{s}},
\qquad
\int_{\R^2\setminus B(0,1)}\frac{\dd x\,\dd y}{(x^2 +
y^2)^{s}} .
\]
推广到 $\R^d$（阈值 $s < d/2$ 与 $s > d/2$）。
\end{exercise}

\begin{exercise}[$\star\star\star$]\label{exo:b3:product:8}
（贝塔--伽马）设 $p, q > 0$，记 $B(p, q) = \int_0^1t^{p-1}(1
- t)^{q-1}\dd t$。将 $\Gamma(p)\Gamma(q)$ 写成二重积分，
作代换 $(x, y) = (uv,\, u(1 - v))$，它将开象限
微分\omterm{def:b3:topology:continuity}{同胚}地映为 $\intoo0\infty\times\intoo01$，
其雅可比行列式 $= u$，并
得出\index{贝塔函数}
\[
B(p, q) = \frac{\Gamma(p)\,\Gamma(q)}{\Gamma(p + q)} .
\]
推出 $\int_0^{\pi/2}\sin^{2p-1}\theta\cos^{2q-1}\theta\,
\dd\theta = \frac12B(p,q)$，并求沃利斯
积分 $W_n = \int_0^{\pi/2}\sin^n$ 的值。
\end{exercise}

\begin{exercise}[$\star\star$]\label{exo:b3:product:9}
（转移公式）设 $T \colon (X, \mathcal A, \mu) \to (Y,
\mathcal B)$ 为可测映射，记 $T_*\mu(B) = \mu(T^{-1}(B))$
为其\emph{像测度}\index{像测度}。
证明：对每个\omterm{def:b3:lebesgue:measurable}{可测函数} $g \geq 0$（定义在 $Y$ 上）有
\[
\int_Y g\,\dd(T_*\mu) = \int_X g\circ T\,\dd\mu
\]
（标准机器）。然后与
\cref{thm:b3:product:linearchange}作比较：变量替换公式
比抽象的转移公式多携带了哪些额外信息？
\emph{（转移公式从不
确定 $T_*\mu$；变量替换定理则计算
$T_*\lambda_d$ 的显式形式，即一个密度\omterm{def:b3:measure:measure}{测度}。）}
\end{exercise}

\begin{exercise}[$\star\star\star$]\label{exo:b3:product:10}
（高斯矩）利用\omterm{ex:b3:product:polar}{极坐标}与富比尼定理，
对 $\R^d$ 上的标准高斯权计算下列积分：
\[
\int_{\R^d}\norm x_2^2\;\eu^{-\norm x_2^2}\,\dd x
\qquad\text{和}\qquad
\int_{\R^d}x_1^2\,\eu^{-\norm x^2_2}\,\dd x,
\]
检验两者的一致性（$\norm x^2 = \sum x_i^2$），并
推出\omterm{def:b3:measure:measure}{测度}
$\pi^{-d/2}\eu^{-\norm x^2}\dd x$ 的二阶矩。
\end{exercise}

\begin{exercise}[$\star\star$]\label{exo:b3:product:11}
（图像与下图）设 $f \colon \R^d \to \intco0\infty$
为\omterm{def:b3:lebesgue:measurable}{可测函数}。
(a) 证明\emph{下图} $H = \{(x, y) \in
\R^d\times\R : 0 < y < f(x)\}$ 是
$\R^{d+1}$ 中的可测集，且
\[
\lambda_{d+1}(H) = \int_{\R^d}f\,\dd\lambda_d :
\]
“积分即图像下方的面积”，终于
得到一条定理。\emph{（截面；托内利。）}
(b) 证明图像 $\{(x, f(x)) : x \in \R^d\}$ 是
$\R^{d+1}$ 中的一个零测集。
(c) 由此给出一个两行的证明，说明球面 $S^{d-1}$ 是
$\R^d$ 中的勒贝格零测集。
\end{exercise}

\begin{exercise}[$\star\star$]\label{exo:b3:product:12}
（一个著名的二重积分）利用几何级数，
并在 $\intoo01^2$ 上使用托内利定理，证明
\[
\int_0^1\!\!\int_0^1\frac{\dd x\,\dd y}{1 - xy}
= \sum_{n\geq1}\frac1{n^2} = \zeta(2),
\qquad
\int_0^1\!\!\int_0^1\frac{\dd x\,\dd y}{1 + xy}
= \sum_{n\geq1}\frac{(-1)^{n-1}}{n^2} = \frac{\zeta(2)}2 .
\]
（第二个级数恒等式：将指标按奇偶分开。）
已知 $\zeta(2) = \frac{\pi^2}6$（见\cref{ch:b3:spectral}），两个看似普通的积分
分别等于 $\frac{\pi^2}6$ 和
$\frac{\pi^2}{12}$；托内利的正性
假设究竟在何处起作用？
\end{exercise}

\section{问题：斯特林公式}

\begin{problem}[{周末问题——$n! \sim
\sqrt{2\pi n}\,(n/\eu)^n$，由控制收敛定理}]
\label{pb:b3:product:1}
斯特林公式支配着本书中每一个渐近计数
——球体体积、二项式系数与中心极限定理的局部形式。
我们从 $\Gamma$ 积分出发证明它，
（见\cref{ex:b3:lebesgue:gamma}），用\emph{拉普拉斯方法}，
以其最简洁的控制收敛形式完成证明，
随后收获诸多应用。\index{斯特林公式}

\medskip
\noindent\textbf{第一部分——公式。}设 $t > 0$，
$\Gamma(t + 1) = \int_0^\infty x^{t}\eu^{-x}\dd x$。
\begin{enumerate}
  \item 作代换 $x = t + \sqrt t\,u$，并证明
        \[
        \frac{\Gamma(t+1)}{t^{t}\eu^{-t}\sqrt t}
        = \int_{-\sqrt t}^{+\infty}
        \exp\Bigl(t\ln\Bigl(1 + \frac u{\sqrt
        t}\Bigr) - \sqrt t\,u\Bigr)\,\dd u
        \;=\;\int_\R g_t(u)\,\dd u,
        \]
        其中 $g_t(u) = \exp\bigl(t\ln(1 + u/\sqrt t) -
        \sqrt t\,u\bigr)\mathbf 1_{u > -\sqrt t}$。
  \item 证明逐点极限：对每个固定的 $u$，$g_t(u)
        \to \eu^{-u^2/2}$，当 $t \to +\infty$ 时 \emph{（将
        $\ln(1 + h)$ 展开到二阶）}。
  \item 控制。设 $\varphi(h) = \ln(1 + h) - h$，于是
        $g_t(u) = \exp\bigl(t\,\varphi(u/\sqrt t)\bigr)$，
        对 $u > -\sqrt t$ 成立，证明如下两个不等式：
        \[
        \varphi(h) \leq -\frac{h^2}4 \quad (-1 < h \leq 1),
        \qquad
        \varphi(h) \leq -c\,h \quad (h \geq 1),\ \ c = 1 -
        \ln 2 > 0
        \]
        \emph{（研究 $\varphi(h) + \frac{h^2}4$ 与 $\varphi(h)
        + ch$：求导数，并检验各区间上的符号
        ）}。由此推得，对 $t \geq 1$：
        \[
        g_t(u) \leq \eu^{-u^2/4}\ \ (\abs u \leq \sqrt t),
        \qquad
        g_t(u) \leq \eu^{-cu}\ \ (u \geq \sqrt t),
        \]
        于是 $g_t(u) \leq \eu^{-u^2/4} +
        \eu^{-cu}\,\mathbf 1_{u > 0}$，即一个\omterm{def:b3:lebesgue:l1}{可积}的控制函数，
        它与 $t \geq 1$ 无关。
  \item 由控制收敛定理与
        高斯积分（见\cref{ex:b3:product:polar}），得出：
        \[
        \Gamma(t + 1) \;\sim\;
        \sqrt{2\pi t}\;\Bigl(\frac t\eu\Bigr)^{t}
        \qquad (t \to +\infty),
        \]
        特别地，$n! \sim \sqrt{2\pi
        n}\,(n/\eu)^n$。
\end{enumerate}

\medskip
\noindent\textbf{第二部分——应用。}
\begin{enumerate}[resume]
  \item （沃利斯）由\cref{exo:b3:product:8}，$W_{2n} =
        \frac\pi2\binom{2n}n4^{-n}$ 型公式，
        由斯特林公式推出 $\binom{2n}{n} \sim \frac{4^n}{\sqrt{\pi n}}$，
        并与递推关系 $W_{n} =
        \frac{n-1}nW_{n-2}$ 相核对。
  \item （球体体积的坍缩）证明
        \[
        v_d = \frac{\pi^{d/2}}{\Gamma(\frac d2 + 1)}
        \;\sim\; \frac{1}{\sqrt{\pi d}}
        \Bigl(\frac{2\pi\eu}{d}\Bigr)^{d/2},
        \]
        于是 $v_d \to 0$ 快于任何几何数列；
        求使 $v_d$ 取最大值的维数（数值计算：$d =
        5$）。
  \item （二项分布的集中——预告
        \cref{ch:b3:clt}）利用斯特林公式证明局部估计，
        即当 $k = n/2 + s\sqrt n/2$，$s$ 固定，
        $n$ 为偶数时：
        \[
        2^{-n}\binom{n}{k} \;\sim\;
        \sqrt{\frac{2}{\pi n}}\;\eu^{-s^2/2},
        \]
        离散高斯轮廓：棣莫弗--拉普拉斯定理的
        雏形。
  \item 第一部分的证明究竟在何处用到了：(i) 单调收敛定理，
        还是控制收敛定理；(ii) 高斯积分；(iii) \omterm{def:b3:measure:lebesgueouter}{勒贝格测度}的不变
        性质？每项各用一句话作答。
\end{enumerate}

\medskip
\noindent\textbf{第三部分——误差项：带界的
斯特林公式。}设 $d_n = \ln n! - \bigl(n + \tfrac12\bigr)\ln n + n -
\ln\sqrt{2\pi}$，于是由第一部分得 $d_n \to 0$。
\begin{enumerate}[resume]
  \item 证明 $d_n - d_{n+1} = \bigl(n +
        \tfrac12\bigr)\ln\bigl(1 + \tfrac1n\bigr) - 1$。
  \item 取 $t = \frac1{2n+1}$，验证 $\frac{n+1}n =
        \frac{1+t}{1-t}$，并展开：
        \[
        d_n - d_{n+1} = \frac{t^2}3 + \frac{t^4}5 +
        \frac{t^6}7 + \cdots,
        \]
        由此推出如下双侧估计
        \[
        \frac1{3(2n+1)^2} \;<\; d_n - d_{n+1} \;<\;
        \frac1{12n} - \frac1{12(n+1)} .
        \]
  \item 裂项求和（利用 $d_m \to 0$），并验证
        漂亮的代数不等式 $\frac1{3(2m+1)^2} >
        \frac1{12m+1} - \frac1{12(m+1)+1}$，对 $m \geq 1$ 成立；
        从而得到经典的夹逼估计：
        \[
        \sqrt{2\pi n}\Bigl(\frac n\eu\Bigr)^n
        \eu^{1/(12n+1)} \;<\; n! \;<\;
        \sqrt{2\pi n}\Bigl(\frac n\eu\Bigr)^n\eu^{1/(12n)} .
        \]
  \item 两个推论：(a) 斯特林公式的相对误差
        $< 10^{-6}$，一旦 $n \geq 83\,334$；
        (b) 利用上述夹逼式，手工估计 $100!$ 到四位有效数字
        （$100! \approx 9.3326\cdot
        10^{157}$），并简要感叹渐近公式在
        极有限的 $n$ 处所达到的精度。
\end{enumerate}

\medskip
\noindent\textbf{第四部分——沃利斯路线：不依赖高斯的
斯特林公式。}历史上，常数 $\sqrt{2\pi}$ 来自沃利斯而非高斯；
本部分重新证明斯特林公式，
独立于第一至二部分，从而重新证明
高斯积分。设 $W_n = \int_0^{\pi/2}\sin^n\theta\,
\dd\theta$。
\begin{enumerate}[resume]
  \item 证明 $W_n = \frac{n-1}nW_{n-2}$（用分部
        积分），并给出闭形式
        \[
        W_{2n} = \frac\pi2\binom{2n}n4^{-n}, \qquad
        W_{2n+1} = \frac{4^n}{(2n+1)\binom{2n}n},
        \]
        及恒等式 $W_nW_{n-1} = \frac\pi{2n}$。
  \item 由 $(W_n)$ 的单调性推出
        $W_{2n}/W_{2n+1} \to 1$，进而
        \[
        W_{2n} \sim \frac12\sqrt{\frac\pi n}
        \qquad\text{且}\qquad
        \binom{2n}n4^{-n}\sqrt n \longrightarrow
        \frac1{\sqrt\pi} :
        \]
        沃利斯定理，无需斯特林公式即可得到。
  \item 仅用第三部分的裂项求和（不需要常数的
        值），证明 $e_n = \ln n! - (n +
        \frac12)\ln n + n$ 收敛于某极限 $\ell$；
        等价地，$n! \sim K\,n^{n+1/2}\eu^{-n}$，其中 $K =
        \eu^\ell > 0$ 尚未确定。
  \item 将此渐近式代入 $\binom{2n}n4^{-n}\sqrt n$，
        并利用第 14 题，确定唯一可能的
        值：$K = \sqrt{2\pi}$。整理逻辑：第三至四部分
        合起来给出完整的第二证明，
        即斯特林公式——由此倒用第一部分的换元，
        独立地算出
        $\int_\R\eu^{-u^2/2}\dd u = \sqrt{2\pi}$。两根支柱，
        二者互为支撑。
\end{enumerate}

\medskip
\noindent\textbf{第五部分——最后的红利。}
\begin{enumerate}[resume]
  \item （完整的局部形态）对整数 $\abs j \leq
        K\sqrt n$（$K$ 固定），证明
        \[
        \frac{\binom{2n}{n+j}}{\binom{2n}{n}} =
        \prod_{i=1}^{\abs j}\frac{n - i + 1}{n + i}
        = \exp\Bigl(-\frac{j^2}n +
        O\Bigl(\frac1{\sqrt n}\Bigr)\Bigr),
        \]
        关于 $j$ 一致成立 \emph{（取对数，并利用 $\ln
        \frac{1-x}{1+y} = -(x + y) + O(x^2 + y^2)$）}。这是
        第 7 题的双侧版本，以及
        在\cref{ch:b3:clt}的周末问题中所引用的精确估计。
  \item （泊松分布预览）利用斯特林公式证明
        $\eu^{-n}\dfrac{n^n}{n!} \sim \dfrac1{\sqrt{2\pi n}}$：
        均值 $n$ 很大的泊松分布的众数
        所承载的质量 $\approx (2\pi n)^{-1/2}$，恰与
        中心极限定理的预测一致。
  \item （\omterm{ex:b3:lebesgue:gamma}{伽马函数}之比）对 $a \in \intoo01$，证明
        $\dfrac{\Gamma(n + a)}{\Gamma(n)\,n^a} \to 1$，利用
        关于对数凸性斜率的界，见
        \cref{pb:b3:lebesgue:1}（即其第 14 题），并
        借助函数方程，推广到每个实数 $a > 0$。
        （这正是“$\Gamma(t+1) \sim$
        斯特林”在整数之间的含义。）
  \item （球的体积，再谈）由 $v_d = \pi^{d/2}/\Gamma(\frac d2
        + 1)$：精确列出 $v_1, \dots, v_7$，
        借助 $\frac{v_d}{v_{d-2}} = \frac{2\pi}d$ 验证单峰性
        （$d < 2\pi$ 时递增，之后递减），并
        证明下述引人注目的母函数恒等式
        \[
        \sum_{k\geq0}v_{2k}\,x^{2k} = \eu^{\pi x^2} :
        \]
        所有偶数维单位球的体积被压缩进同一个
        指数函数之中。
  \item （熵的渐近性）对固定的 $\alpha \in \intoo01$
        且 $\alpha n \in \N$，由斯特林公式推出
        \[
        \binom{n}{\alpha n} \;\sim\;
        \frac{\eu^{n\,H(\alpha)}}
        {\sqrt{2\pi\,\alpha(1-\alpha)\,n}},
        \qquad
        H(\alpha) = -\alpha\ln\alpha -
        (1-\alpha)\ln(1-\alpha) :
        \]
        二项式系数的指数增长率
        是\emph{熵} $H$——验证 $\alpha =
        \frac12$ 时恢复第 5 题的结果，并验证 $H(\alpha) <
        \ln2$ 对 $\alpha \neq \frac12$ 成立（从而偏离中心的
        二项式相对于 $2^n$ 指数可忽略）。
  \item （表面积）单位球面 $S^{d-1}$
        的面积为 $s_{d-1} = d\,v_d$（作为
        \cref{exo:b3:forms:6} 在微分形式一章中证明；
        此处将其取为定义）。列出
        $s_0, \dots, s_6$，找出最大者（$d - 1 =
        6$，$s_6 = \frac{16\pi^3}{15} \approx 33.07$），并
        证明 $s_{d-1} \to 0$ 同样快于任何几何速率——
        高维球面按任一欧几里得度量都是
        极其微小的。
  \item （第一修正项）由
        第 11 题的夹逼式推出 $d_n = \frac1{12n} +
        O\bigl(\frac1{n^2}\bigr)$，于是
        \[
        n! = \sqrt{2\pi n}\,\Bigl(\frac
        n\eu\Bigr)^{n}\Bigl(1 + \frac1{12n} +
        O\Bigl(\frac1{n^2}\Bigr)\Bigr).
        \]
        取 $n = 10$ 验算：未修正的公式给出
        $3\,598\,696$（相对误差 $8.3\cdot10^{-3}$），
        修正后的结果为 $3\,628\,685$，对照 $10! =
        3\,628\,800$（相对误差 $3.2\cdot10^{-5}$）——
        级数的一项换来两位半的有效数字。
  \item （$\Gamma$ 函数的中位数）证明
        \[
        \frac{1}{\Gamma(t+1)}
        \int_0^{t} x^{t}\eu^{-x}\,\dd x
        \;\longrightarrow\; \frac12
        \qquad (t \to +\infty) :
        \]
        渐近地说，恰有一半的
        $\Gamma$ 被积函数的质量位于其众数 $x = t$ 之下。
        \emph{（对截断积分施行第一部分的换元；
        第 3 题中的分母已经
        就位。）}
  \item （熵，非渐近形式）对 $\alpha \in
        \intoc0{\frac12}$，证明下列界，它对
        \emph{每个} $n \geq 1$ 都成立：
        \[
        \sum_{k=0}^{\lfloor\alpha n\rfloor}\binom nk
        \;\leq\; \eu^{n\,H(\alpha)} ,
        \]
        比较该和式与 $\sum_k\binom
        nk\lambda^{k-\alpha n}$，其中取倾斜参数 $\lambda =
        \frac{\alpha}{1-\alpha} \leq 1$。验证这一
        $\lambda$ 的选取是最优的，并与
        第 21 题对照：渐近陈述中的指数增长率 $H(\alpha)$
        由一个单行不等式达到，
        且完全不涉及渐近分析。
\end{enumerate}
\end{problem}
